Este trabalho examinou se representações de letras acrescentam informação preditiva a atributos acústicos e metadados na estimação da popularidade Spotify observada. A construção dos dados foi parte central da análise: de um catálogo acessível com 1.159.764 faixas, selecionaram-se 200.000 por amostragem estratificada e recuperaram-se letras aceitas para 77.535. A descrição mostrou que esse domínio difere do catálogo em popularidade e cobertura. A seleção probabilística inicial não autoriza tratar automaticamente o subconjunto textual como representativo da população de origem.

Na comparação Ridge com validação agrupada aninhada comum, o MAE passou de 8,3525 no Baseline a 8,1151 no Combinado, redução de 2,84\%. A configuração com embeddings apresentou a maior parcela desse ganho; o acréscimo posterior dos 21 descritores produziu redução menor, de 0,0134 ponto, inconclusiva na análise auxiliar de incerteza. A sensibilidade com pesos originais apenas na avaliação preservou a ordenação, sem corrigir automaticamente a disponibilidade de letras.

O resultado médio não se reproduziu uniformemente nos subgrupos. Todas as 128 faixas aceitas de popularidade elevada foram subestimadas, e os casos de idioma incerto não apresentaram ganho conjunto. Previsões negativas decorreram de um estimador sem restrição de suporte; a sensibilidade por clipping reduziu ligeiramente o MAE, mas não melhorou a cauda. Assim, a evidência sustenta contribuição textual modesta no procedimento e domínio estudados, sem demonstrar capacidade prospectiva de prever sucessos ou efeito causal das letras.

Como continuidade, são pertinentes a confirmação de metadados temporais da fonte, uma revisão humana planejada da qualidade de correspondência e a avaliação independente de uma representação com prefixo corrigido. Métodos para ampliar o alcance inferencial dependeriam de hipóteses explícitas sobre seleção textual e suporte, além de validação apropriada. Essas possibilidades constituem novos estudos delimitados, sem substituir ou reinterpretar retrospectivamente a comparação aqui apresentada.
